\section{Proxy Validity on Literal Observed Paths}
\label{sec:observation}

Let $P_t$ denote in-domain perplexity and $E_t$ a knowledge-sensitive
evaluation at an observed checkpoint $t$. We examine only the implication
\[
P_{t_2}<P_{t_1}\quad\Longrightarrow\quad
\text{target recovery at }t_2 .
\]
Here ``target recovery'' means closing the large deficit to the intact-base
score on the same evaluation, not merely obtaining a statistically detectable
within-run gain. A path contradicts this improvement-only implication when PPL
falls while a large intact-base target deficit remains. This is deliberately
weaker than evaluating a prospective certificate: no absolute PPL threshold,
base-relative tolerance, plateau rule, or construction-specific calibration was
pre-registered. The analysis cannot identify where any capability resides,
establish why paths differ, or predict behavior after the final checkpoint.
The best observed keep14 PPL remains $1.428\times$ the base and only three late
checkpoints are available, so we do not retrofit a post-hoc threshold table.

This framing motivates three design choices. We preserve within-run checkpoints
rather than infer a universal dynamic from an endpoint. We report multiple
evaluation interfaces because each operationalizes the target differently. We
also show null and alternative-construction operating points to bound simple
explanations, while exposing every unmatched factor rather than calling the
comparisons controlled ablations.
